\subsection{The property (GDF)}

In this No.~we are going to consider locally convex spaces F having the following property (called the `countably closed graph' property): \(^{6}\)

(GDF) If u is a linear mapping of F into a Banach space B such that, in the product space \(F \times B\) , every limit of every convergent sequence of points of the graph \(\Gamma\) of u again belongs to \(\Gamma\) , then u is continuous.

Every Fréchet space has the property (GDF) (TVS, I, §3, No.~3, Cor. 5 of Th. 1). In the Appendix, we shall see other examples of spaces having the property (GDF).

PROPOSITION 11. --- Every Hausdorff locally convex space F with the property (GDF) is barreled.

Let V be a barrel in F, q is gauge, which is a semi-norm on F; let H be the Hausdorff space associated with the space F equipped with the topology defined by this single semi-norm. The completion \(\hat{H}\) of H is a Banach space; let \(\pi\) be the canonical mapping of F into \(\hat{H}\) ; we are going to show that \(\pi\) is continuous (for the original topology on F); this will establish the proposition, because V, the inverse image under \(\pi\) of the unit ball of \(\hat{H}\) , will then be a neighborhood of 0 in F. To establish the continuity of \(\pi\) it will suffice, by virtue of (GDF), to show that the graph of \(\pi\) is closed in \(F \times \hat{H}\) ; in other words, we must see that if \(\mathfrak{F}\) is a filter on F, convergent to \(x \in F\) , and its image \(\pi(\mathfrak{F})\) converges to \(y \in \hat{H}\) , then \(y = \pi(x)\) . Now, every element \(x'\) of the polar \(V^{\circ}\) of V in \(F'\) may be extended in only one way to a continuous linear form on \(\hat{H}\) (again denoted \(x'\) ), and the set of these forms is the unit ball of the dual of \(\hat{H}\) ; it therefore suffices to show that \(\langle y, x' \rangle = \langle \pi(x), x' \rangle\) for every \(x' \in V^{\circ}\) . But this follows from the relations

\[
\langle \mathbf {y}, \mathbf {x} ^ {\prime} \rangle = \lim _ {\mathfrak {F}} \left\langle \pi (\mathbf {z}), \mathbf {x} ^ {\prime} \right\rangle = \lim _ {\mathfrak {F}} \left\langle \mathbf {z}, \mathbf {x} ^ {\prime} \right\rangle = \left\langle \mathbf {x}, \mathbf {x} ^ {\prime} \right\rangle = \left\langle \pi (\mathbf {x}), \mathbf {x} ^ {\prime} \right\rangle .
\]

THEOREM 1. (Gelfand--Dunford)--- Let F be a Hausdorff locally convex space having the property (GDF), \(F_{s}^{\prime}\) its weak dual. For every mapping f of T into \(F_{s}^{\prime}\) that is scalarly essentially \(\mu\) -integrable, the integral \(\int f d\mu\) belongs to \(F^{\prime}\) .

Recall that the dual of \(F_{s}^{\prime}\) is F (TVS, II, §6, No.~2, Prop. 3). For every \(z \in F\) , the numerical function \(\langle z, f \rangle\) is therefore essentially \(\mu\) -integrable; let \(\theta(z)\) be its class in \(L^{1}(\mu)\) . To show that \(\int f d\mu \in F^{\prime}\) , one must establish that the linear form \(z \mapsto \langle z, \int f d\mu \rangle\) is continuous on F; in fact, we are going to prove the following stronger result:

Lemma 2. --- Let f be a mapping of T into \(F_{s}^{\prime}\) , such that for every \(z \in F\) , the numerical function \(\langle z, f \rangle\) belongs to \(\overline{\mathcal{L}^{p}}(\mu) (1 \leqslant p \leqslant +\infty)\) ; let \(\theta(z)\) be the class of this function in \(L^{p}(\mu)\) . Then \(z \mapsto \theta(z)\) is a continuous linear mapping of F into \(L^{p}(\mu)\) .

By virtue of the (GDF) property, it suffices to show that for every sequence \((\mathbf{z}_{n})\) of elements of F converging to z, such that \((\theta(\mathbf{z}_{n}))\) converges to \(u \in \mathrm{L}^{p}(\mu)\) , one has \(u = \theta(\mathbf{z})\) . Now, replacing if necessary the sequence \((\mathbf{z}_{n})\)

by a subsequence, one can suppose that the sequence of functions \(\langle \mathbf{z}_n,\mathbf{f}\rangle\) converges locally almost everywhere to a function \(h\in \overline{\mathcal{L}^p} (\mu)\) , with class \(u\) in \(\mathrm{L}^p (\mu)\) (Ch. IV, §3, No.~4, Th. 3 and Ch. V, §1, No.~3). Since, by hypothesis, for every \(t\in \mathrm{T}\) the sequence \((\langle \mathbf{z}_n,\mathbf{f}(t)\rangle)\) converges to \(\langle \mathbf{z},\mathbf{f}(t)\rangle\) , we have \(h(t) = \langle \mathbf{z},\mathbf{f}(t)\rangle\) locally almost everywhere, consequently \(u = \theta (\mathbf{z})\) .

COROLLARY 1. --- Let \(G_{i}\) ( \(1 \leqslant i \leqslant n\) ) be n Hausdorff locally convex spaces having the property (GDF), and let F be the space of separately continuous multilinear forms on \(\prod_{i=1}^{n} G_{i}\) , equipped with the topology of pointwise convergence. For every mapping f of T into F that is scalarly essentially \(\mu\) -integrable, one has \(\int f d\mu \in F\) .

The space F is in duality with the tensor product \(\bigotimes_{i=1}^{n}G_{i}\) , and the topology of pointwise convergence on F is none other than the topology \(\sigma(F,\bigotimes_{i=1}^{n}G_{i})\) . The algebraic dual \(F^{\prime*}\) is therefore the space of all multilinear forms on \(\prod_{i=1}^{n}G_{i}\) . Let \(\mathbf{z}=(\mathbf{z}_{1},\ldots,\mathbf{z}_{n})\) be an element of \(\prod_{i=1}^{n}G_{i}\) ; for every multilinear form \(u\in F^{\prime*}\) , the mapping \(\mathbf{x}\mapsto\mathbf{u}(\mathbf{z}_{1},\ldots,\mathbf{z}_{i-1},\mathbf{x},\mathbf{z}_{i+1},\ldots,\mathbf{z}_{n})\) is a linear form on \(G_{i}\) , which we shall denote by \(\lambda_{i}(\mathbf{z})(\mathbf{u})\) ; we thus obtain a linear mapping \(\lambda_{i}(\mathbf{z})\) of \(F^{\prime*}\) into the algebraic dual \(G_{i}^{*}\) of \(G_{i}\) , continuous for the topologies \(\sigma(F^{\prime*},\bigotimes_{i=1}^{n}G_{i})\) and \(\sigma(G_{i}^{*},G_{i})\) . To say that \(u\in F\) means that for every index i and every \(z\in\Pi_{i=1}^{n}G_{i}\) , one has \(\lambda_{i}(\mathbf{z})(\mathbf{u})\in G_{i}^{\prime}\) . Now, by Prop. 1 of No.~1, the mapping \(\lambda_{i}(\mathbf{z})\circ\mathbf{f}\) is a scalarly essentially \(\mu\) -integrable mapping of T into \(G_{i}^{\prime}\) equipped with the topology \(\sigma(G_{i}^{\prime},G_{i})\) , and

\[
\int \left(\lambda_ {i} (\mathbf {z}) \circ \mathbf {f}\right) d \mu = \lambda_ {i} (\mathbf {z}) \left(\int \mathbf {f} d \mu\right).
\]

By Th. 1, \(\int (\lambda_i(\mathbf{z})\circ \mathbf{f})d\mu \in \mathrm{G}_i^\prime\) for \(1\leqslant i\leqslant n\) , therefore \(\int \mathbf{f}d\mu \in \mathbf{F}\)

COROLLARY 2. --- Let G be a Hausdorff locally convex space having the property (GDF), and H a semi-reflexive space whose strong dual \(H_{b}^{\prime}\) has the property (GDF) (cf.~App., No.~2, Prop. 3). Let F be the space \(\mathcal{L}_{s}(G;H)\) ; for every mapping U of T into F that is scalarly essentially \(\mu\) -integrable, the integral \(\int U d\mu\) belongs to F.

Since G is barreled (Prop. 11), \(\mathcal{L}(\mathrm{G};\mathrm{H}) = \mathcal{L}(\mathrm{G}_{\sigma};\mathrm{H}_{\sigma})\) (TVS, IV, §1, No.~3, Prop. 7); moreover, one can replace \(\mathrm{F} = \mathcal{L}_s(\mathrm{G};\mathrm{H})\) by the space \(\mathcal{L}_s(\mathrm{G}_{\sigma};\mathrm{H}_{\sigma})\) , since the two spaces have the same dual \(\mathrm{G} \otimes \mathrm{H}'\) (TVS, II, §6, No.~2, Prop. 3, and the first paragraph of No.~3 above). If, for every \(u \in \mathcal{L}(\mathrm{G};\mathrm{H}) = \mathcal{L}(\mathrm{G}_{\sigma};\mathrm{H}_{\sigma})\) , one sets \(\widetilde{u}(\mathbf{x},\mathbf{y}') = \langle u(\mathbf{x},\mathbf{y}') \rangle\) (for \(\mathbf{x} \in \mathbf{G}\) , \(y' \in H')\) , the linear mapping \(u \mapsto \widetilde{u}\) is a bijection of F onto the space \(F_1\) of separately continuous bilinear forms on \(G_\sigma \times H_s'\) , where \(H_s'\) denotes the dual \(H'\) equipped with the weak topology \(\sigma(H', H)\) (App., No.~1); moreover, this mapping is an isomorphism of \(\mathcal{L}_s(G_\sigma; H_\sigma)\) onto \(F_1\) equipped with the topology of pointwise convergence (loc. cit.). But since by hypothesis H is the dual of \(H_b'\) , \(F_1\) is also the space of separately continuous bilinear forms on \(G \times H_b'\) . Cor. 2 therefore follows from Cor. 1.

Note that Cor. 2 is applicable in particular when G is Banach space and H is a reflexive Banach space.

